\section{Array: Membuat, Mengakses Indeks, dan Metode Dasar}

Array adalah struktur data yang menyimpan kumpulan nilai dalam urutan terindeks. Indeks dimulai dari 0. Array di JavaScript bersifat dinamis---ukurannya dapat berubah dan isinya dapat berisi tipe data campuran (meskipun tidak direkomendasikan). Array dibuat dengan literal \code{[]} atau \code{new Array()} (literal lebih direkomendasikan) \cite{jsInfo}.

\begin{table}[htbp]
\centering
\begin{tabular}{|P{3.5cm}|P{4cm}|P{4.5cm}|}
\hline
\textbf{Metode} & \textbf{Fungsi} & \textbf{Contoh} \\
\hline
\code{.push(item)} & Tambah di akhir & \code{arr.push(4)} \\
\hline
\code{.pop()} & Hapus dari akhir & \code{arr.pop()} \\
\hline
\code{.unshift(item)} & Tambah di awal & \code{arr.unshift(0)} \\
\hline
\code{.shift()} & Hapus dari awal & \code{arr.shift()} \\
\hline
\code{.indexOf(item)} & Cari posisi & \code{arr.indexOf(3)} \\
\hline
\code{.includes(item)} & Cek keberadaan & \code{arr.includes(2)} \\
\hline
\code{.slice(s, e)} & Salin bagian & \code{arr.slice(1, 3)} \\
\hline
\code{.splice(i, n)} & Hapus/sisip di posisi & \code{arr.splice(1, 2)} \\
\hline
\code{.forEach(fn)} & Iterasi setiap elemen & \code{arr.forEach(x => ...)} \\
\hline
\code{.map(fn)} & Transformasi ke array baru & \code{arr.map(x => x * 2)} \\
\hline
\code{.filter(fn)} & Filter ke array baru & \code{arr.filter(x => x > 2)} \\
\hline
\end{tabular}
\caption{Metode array JavaScript yang paling sering digunakan}
\end{table}

\begin{webcode}[language=JavaScript, caption={Operasi array dasar dan fungsional}]
const numbers = [10, 20, 30, 40, 50];

// Akses dan modifikasi
console.log(numbers[0]);   // 10
console.log(numbers.length); // 5
numbers.push(60);  // [10, 20, 30, 40, 50, 60]

// Metode fungsional (menghasilkan array baru)
const doubled = numbers.map(n => n * 2);
const big = numbers.filter(n => n > 30);
const total = numbers.reduce((sum, n) => sum + n, 0);
console.log(doubled); // [20, 40, 60, 80, 100, 120]
console.log(total);   // 210
\end{webcode}

\begin{konsep}
\textbf{Destructuring Array (ES6).} Sintaks \code{const [a, b, ...rest] = arr;} mengekstrak elemen array ke variabel individual. Rest operator (\code{...rest}) menampung sisa elemen. Destructuring juga bekerja untuk swap nilai: \code{[a, b] = [b, a];}. Ini membuat kode lebih ringkas dan ekspresif \cite{mdnJsGuide}.
\end{konsep}

